% s2: begin with the geometry. ASCII only.
\section*{2. Begin with the geometry}

I start from the geometry the theory already lives on, a principal fibre bundle, equipped with a connection and a
space-time metric. The \emph{base} is space-time (or a spatial slice in the Hamiltonian picture).
Locally its fibres are copies of the \emph{structure group} $SU(2)$, with no preferred group identity in each fibre.
Its double cover of $SO(3)$ concerns internal colour rotations, not space-time spin. Its centre $\{\pm1\}$ permits
half-integer colour representations and distinguishes fundamental from adjoint loops; it does not alone prove
confinement~\cite{tHooft78}. A \emph{connection} $A$ is the gauge field, a rule for
carrying the fibre along a path; around a closed loop $C$ it returns a \emph{holonomy} $\mathrm{Hol}_C(A)\in SU(2)$, and
its infinitesimal holonomy is the curvature $F$. A \emph{gauge transformation} relabels each fibre, so the physical
configurations are connections modulo relabelling, the space $\mathcal A/\mathcal G$. That space is curved and folded,
and no single choice of labels covers it (the Gribov problem~\cite{Gribov78,Singer78}). It has \emph{topology}: winding
classes of large gauge transformations $\pi_3(SU(2))=\mathbb Z$. Chern--Simons shifts by that integer; on general
connections it is continuous, not an integer-valued topological label~\cite{ChernSimons74}. Vacuum representatives are joined by
instantons~\cite{BPST75,JackiwRebbi76,CallanDashenGross76}. The quantum theory adds a measure, a spectrum, and a sorting of
modes into slow and fast.

\textbf{Three things working together.} I do not ask topology to assign an energy by itself. It describes the global
connections and identifications of the space of physical fields. The spatial metric and the Yang--Mills action give
that space its ruler and its magnetic-energy landscape. The kinetic term measures the energy cost of variation of a
wavefunctional across field configurations. Quantum dynamics combines them, with the gauge constraints, to determine
the vacuum and the excitations~\cite{KogutSusskind75,vanBaal00}. These are aspects of the same theory, not three extra
forces. The same topology can carry different geometries and different spectra. I want to follow what changes in
the field, not only the number attached to that change. Van Baal describes how wavefunctionals spreading across
identified boundaries of configuration space become sensitive to its topology~\cite{vanBaal92Topology}; that is a
physical connection between the global shape of the space and its low energies.

A connection component at a point depends on the gauge; local invariant quantities such as $\mathrm{tr}\,F^2$ do not.
Loops give another set of physical readings:
$W_C=\mathrm{tr}_\rho\,\mathrm{Hol}_C(A)$ is unchanged by relabelling, because the trace does not see conjugation, and
the complete compatible loop data reconstruct a smooth connection up to gauge under the hypotheses of
Giles' theorem~\cite{Giles81}. A finite set of measured loops is not that complete data. Wilson's lattice theory~\cite{Wilson74} is this
bundle sampled: the base at the vertices of a grid of spacing $a$, a group element $U_\ell$ on each edge, and the action
read from the holonomy $U_p$ of each elementary square (plaquette).

\begin{claim}{PROVED}{The lattice measure lives on connections modulo gauge}\label{b:wilson}
$d\mu_{a,L}=Z^{-1}\exp\bigl(-\beta\sum_p(1-\tfrac12\mathrm{tr}\,U_p)\bigr)\prod_\ell dU_\ell$, $\beta=4/g^2$, is gauge
invariant and descends to a probability measure on lattice connections modulo gauge. No gauge fixing is needed for
these integrals; this avoids, rather than solves, the Gribov gauge-fixing obstruction. Every Wilson loop is gauge invariant.
\evidence{Traces of holonomies and an invariant Haar measure~\cite{Wilson74}; checked in my engine (\S3).}
\lines{1}
\end{claim}

\begin{claim}{PROVED (CITED)}{Reflection positivity and the transfer matrix}\label{b:rp}
The Wilson measure is reflection positive, so there are a Hilbert space, a vacuum and a transfer matrix $T=e^{-aH}$ with
$H\ge0$.
\evidence{Osterwalder and Seiler~\cite{OsterwalderSeiler78}, in the sense of~\cite{OS73,OS75}.}
\lines{2}
\end{claim}

Two choices shape the rest. The size of the room is a coordinate of a separate parameter space, not a nuisance, so the gap is a function
over the family of rooms (\S7). The physical base is not being identified with this parameter space.
The rooms are spheres throughout; I retain homogeneous fields as a slow sector with topological information,
without assuming they exhaust the low-energy physical states. Infinite space is the proposed end of the family (\S8). First, the lattice
already gives the gap a meaning at strong coupling, where it can be read off by hand.
